\documentclass[12pt]{article}
\usepackage[utf8]{inputenc}
\usepackage[T1]{fontenc}
\usepackage[margin=1in]{geometry}
\usepackage{setspace}
\usepackage{xcolor}
\usepackage[colorlinks=true, urlcolor=blue!60!black, linkcolor=blue!60!black]{hyperref}
\setlength{\parindent}{2em}
\setlength{\parskip}{1.25em}
\renewcommand{\baselinestretch}{1.0}
\title{\textit{Sveen v. Melin}, 584 U.S. 811 (2018)}
\author{Case Brief}
\date{}
\begin{document}
\maketitle

\section{Case Metadata}
\begin{itemize}
  \item \textbf{Case Name:} \textit{Sveen v. Melin}, 584 U.S. 811 (2018) (138 S. Ct. 1815; Docket No. 16-1432).
  \item \textbf{Dates:} Argued March 19, 2018; Decided June 11, 2018.
  \item \textbf{Vote and Alignment:} 8--1, reversing and remanding the Eighth Circuit (\textit{Metropolitan Life Ins. Co. v. Melin}, 853 F.3d 410 (8th Cir. 2017)).
  \item \textbf{Majority Opinion:} Justice Kagan, joined by Chief Justice Roberts and Justices Kennedy, Thomas, Ginsburg, Breyer, Alito, and Sotomayor.
  \item \textbf{Dissenting Opinion:} Justice Gorsuch (sole dissenter).
  \item \textbf{Constitutional Provision:} Contracts Clause, U.S. Const. art. I, \S~10, cl. 1 (``No State shall \ldots{} pass any \ldots{} Law impairing the Obligation of Contracts'').
  \item \textbf{Primary Opinion Sources:} \href{https://www.supremecourt.gov/opinions/17pdf/584us2r49_74h3.pdf}{SupremeCourt.gov (Official Slip/Proof Opinion)}, \href{https://www.law.cornell.edu/supremecourt/text/16-1432}{Cornell LII}, and \href{https://supreme.justia.com/cases/federal/us/584/16-1432/}{Justia U.S. Supreme Court Center}.
\end{itemize}

\section{Facts and Procedural History}
In 1997, Mark Sveen purchased a revocable life-insurance policy from Metropolitan Life Insurance Company (MetLife) and married respondent Kaye Melin later that year. In 1998, Mark named Melin as the policy's primary beneficiary and designated his two adult children from a prior marriage, petitioners Ashley Sveen and Antone Sveen, as contingent beneficiaries.

In 2002, Minnesota amended its probate code to extend its revocation-on-divorce rule from wills to ``will substitutes'' such as life-insurance policies. Codified at Minn. Stat. \S~524.2-804, subd. 1, the statute provides that ``the dissolution or annulment of a marriage revokes any revocable \ldots{} beneficiary designation \ldots{} made by an individual to the individual's former spouse,'' treating the ex-spouse as having predeceased the policyholder unless the policyholder affirmatively re-designates the ex-spouse after divorce or the divorce decree provides otherwise.

Mark Sveen and Kaye Melin divorced in 2007. Their divorce decree did not mention the MetLife policy, and Mark never altered his beneficiary designation. When Mark died in 2011, both Melin and the Sveen children claimed the insurance proceeds. MetLife filed a federal interpleader action in the U.S. District Court for the District of Minnesota to determine the rightful payee.

The District Court granted summary judgment for the Sveen children, holding that retroactive application of Minnesota's 2002 statute revoked Melin's designation and did not violate the Contracts Clause. The U.S. Court of Appeals for the Eighth Circuit reversed (853 F.3d 410 (2017)), holding that retroactively applying a revocation-on-divorce statute to a pre-existing life-insurance designation disrupted the policyholder's expectations and violated the Contracts Clause. To resolve a circuit and state-court split---contrasting the Eighth Circuit and Pennsylvania Supreme Court with the Ninth Circuit, Tenth Circuit, and Colorado Supreme Court---the Supreme Court granted certiorari on December 8, 2017.

\section{Issue(s) Presented}
Does the retroactive application of Minnesota's revocation-on-divorce statute (Minn. Stat. \S~524.2-804, subd. 1) to a life-insurance beneficiary designation made before the statute's enactment violate the Contracts Clause of Article I, Section 10, Clause 1?

\section{Holding}
No. Applying the Court's established two-step Contracts Clause framework, retroactive application of Minnesota's revocation-on-divorce statute does not violate the Contracts Clause because the law fails to cross the step-one threshold of ``substantial impairment'' of a pre-existing contractual relationship.

\section{Reasoning}
Writing for an eight-Justice majority, Justice Kagan began by reaffirming that while the Contracts Clause restricts state disruption of private contracts, not every state law affecting pre-existing agreements violates the Clause (\textit{El Paso v. Simmons}, 379 U.S. 497, 506--507 (1965)). The Court governs Contracts Clause challenges under a two-step inquiry:
\begin{enumerate}
  \item \textbf{Threshold Substantial Impairment:} Whether the state law has ``operated as a substantial impairment of a contractual relationship'' (\textit{Allied Structural Steel Co. v. Spannaus}, 438 U.S. 234, 244 (1978)), evaluated by examining the extent to which the law (a) undermines the contractual bargain, (b) interferes with a party's reasonable expectations, and (c) prevents the party from safeguarding or reinstating his or her rights.
  \item \textbf{Means-Ends Scrutiny:} Only if a substantial impairment exists does the Court ask whether the state law is drawn in an ``appropriate'' and ``reasonable'' way to advance ``a significant and legitimate public purpose'' (\textit{Energy Reserves Group, Inc. v. Kansas Power \& Light Co.}, 459 U.S. 400, 411--412 (1983)).
\end{enumerate}

The Court resolved the case entirely at step one---declining in footnote 2 to revisit Melin's invitation to abandon the modern two-step test in favor of an originalist standard---because three combined features of Minnesota's statute negated any substantial impairment:
\begin{itemize}
  \item \textbf{Alignment with Policyholder Intent:} Revocation-on-divorce statutes rest on the empirical presumption that ``the average Joe does not want his ex inheriting what he leaves behind.'' Because divorced policyholders routinely forget to update nonprobate designations, the law is designed to approximate probable intent and support, rather than subvert, the contractual scheme.
  \item \textbf{No Disruption of Reasonable Expectations:} At the time of contracting, no policyholder could reasonably expect a spousal beneficiary designation to remain insulated from divorce. State divorce courts have long exercised broad equitable authority to divide marital property and reallocate life-insurance policies upon dissolution; the statute merely automates what a divorce court could always have ordered.
  \item \textbf{Minimal Paperwork Burden (Reversible Default Rule):} The statute establishes a mere default rule that a policyholder who wishes to keep an ex-spouse as beneficiary can override ``with the stroke of a pen'' by submitting a standard change-of-beneficiary form to the insurer.
\end{itemize}

Justice Kagan anchored this third rationale in a two-century line of Contracts Clause precedents upholding retroactive recording, notice, and filing laws---including \textit{Jackson v. Lamphire}, 28 U.S. (3 Pet.) 280 (1830), \textit{Curtis v. Whitney}, 80 U.S. (13 Wall.) 68 (1872), \textit{Vance v. Vance}, 108 U.S. 514 (1883), \textit{Gilfillan v. Union Canal Co. of Pa.}, 109 U.S. 401 (1883), \textit{Conley v. Barton}, 260 U.S. 677 (1923), and \textit{Texaco, Inc. v. Short}, 454 U.S. 516 (1982). Those laws extinguished entire property or debt rights for noncompliance with new filing requirements, whereas Minnesota's statute merely redirects insurance proceeds to contingent beneficiaries. Finally, the Court rejected Melin's proposed distinction between laws that alter substantive contract ``rights'' and laws that regulate procedural ``remedies,'' citing \textit{Texaco} for the rule that when eliminating a remedy and extinguishing a right produce identical practical consequences, the Contracts Clause analysis is the same.

\section{Concurring and Dissenting Opinions}
\textbf{Justice Gorsuch (Dissenting):} Writing alone, Justice Gorsuch challenged both the doctrinal foundation of the modern balancing test and its application to Minnesota's statute. First, drawing on the text of Article I, Section 10 and Framing-era history (\textit{Sturges v. Crowninshield}, 17 U.S. (4 Wheat.) 122 (1819)), he argued that the Contracts Clause categorically forbids \textit{any} retroactive legislative impairment of existing contracts---not merely ``substantial'' ones---and warned that critics of the post-\textit{Blaisdell} balancing framework ``deserve a thoughtful reply, if not in this case then in another.''

Second, even under the majority's two-step test, Justice Gorsuch contended that Minnesota's law failed both prongs. At step one, naming the beneficiary is the ``whole point'' of a life-insurance contract; retroactively stripping a named beneficiary alters a core contractual term rather than merely adding a procedural recording step. He dismissed the majority's three justifications: (1) legislatures cannot reliably guess every divorcing spouse's intent (many ex-spouses retain policies for child support or amicable financial security); (2) a judicial divorce decree dividing specific property is fundamentally different from a blanket legislative rewrite of contracts that a divorce court left untouched; and (3) a default rule that catches inattentive policyholders unawares still destroys the contract they actually formed. At step two, he argued that Minnesota had less intrusive, more moderate alternatives---such as requiring divorce courts to notify spouses to review beneficiary forms---without retroactively rewriting pre-existing contracts.

\section{Key Quotations}
\begin{quote}
``In answering that question, the Court has considered the extent to which the law undermines the contractual bargain, interferes with a party's reasonable expectations, and prevents the party from safeguarding or reinstating his rights.'' \\
--- Justice Kagan, for the Court, 584 U.S. at 819 (\href{https://www.law.cornell.edu/supremecourt/text/16-1432}{Cornell LII}; \href{https://supreme.justia.com/cases/federal/us/584/16-1432/}{Justia}).
\end{quote}

\begin{quote}
``First, the statute is designed to reflect a policyholder's intent---and so to support, rather than impair, the contractual scheme. Second, the law is unlikely to disturb any policyholder's expectations because it does no more than a divorce court could always have done. And third, the statute supplies a mere default rule, which the policyholder can undo in a moment.'' \\
--- Justice Kagan, for the Court, 584 U.S. at 820 (\href{https://www.law.cornell.edu/supremecourt/text/16-1432}{Cornell LII}; \href{https://supreme.justia.com/cases/federal/us/584/16-1432/}{Justia}).
\end{quote}

\begin{quote}
``In cases going back to the 1800s, this Court has held that laws imposing such minimal paperwork burdens do not violate the Contracts Clause.'' \\
--- Justice Kagan, for the Court, 584 U.S. at 822 (\href{https://www.law.cornell.edu/supremecourt/text/16-1432}{Cornell LII}).
\end{quote}

\begin{quote}
``[The Framers] took the view that treating existing contracts as `inviolable' would benefit society by ensuring that all persons could count on the ability to enforce promises lawfully made to them---even if they or their agreements later prove unpopular with some passing majority.'' \\
--- Justice Gorsuch, dissenting, 584 U.S. at 828 (\href{https://www.supremecourt.gov/opinions/17pdf/584us2r49_74h3.pdf}{SupremeCourt.gov}; \href{https://www.law.cornell.edu/constitution-conan/article-1/section-10/clause-1/historical-background-on-the-contract-clause}{Cornell LII ConAn}).
\end{quote}

\section{Doctrinal Significance and Modern Legacy (Through 2026)}
\textit{Sveen v. Melin} was the Supreme Court's first major Contracts Clause decision in decades, confirming that the modern two-step framework descended from \textit{Home Building \& Loan Ass'n v. Blaisdell}, 290 U.S. 398 (1934), \textit{Allied Structural Steel}, and \textit{Energy Reserves} remains firmly entrenched---even on an increasingly originalist Court. By commanding an 8--1 majority (including Chief Justice Roberts and Justices Thomas and Alito), Justice Kagan signaled that the Court was not ready to revive Lochner-era or Marshall-Court strictness for state economic or domestic-relations default rules.

Yet \textit{Sveen}'s refined three-factor formulation of ``substantial impairment''---and Justice Gorsuch's originalist dissent---became the primary battleground for a wave of COVID-19 and housing-regulation Contracts Clause litigation between 2020 and 2026:
\begin{itemize}
  \item \textbf{Pandemic Eviction Moratoria and Guaranty Laws:} Courts split sharply when applying \textit{Sveen}'s two-step test to pandemic housing measures. In \href{https://law.justia.com/cases/federal/appellate-courts/ca9/20-56251/20-56251-2021-08-25.html}{\textit{Apartment Ass'n of Los Angeles County v. City of Los Angeles}, 10 F.4th 905 (9th Cir. 2021)}, the Ninth Circuit assumed substantial impairment under \textit{Sveen} but upheld Los Angeles's temporary eviction moratorium at step two as a reasonable response to a public-health emergency; similarly, in \href{https://www.gmsr.com/wp-content/uploads/2025/04/Iten-v.-County-of-LA-24-2974.pdf}{\textit{Iten v. County of Los Angeles}, No. 24-2974 (9th Cir. 2025)}, the Ninth Circuit rejected a commercial landlord's challenge at \textit{Sveen} step one. By contrast, in \href{https://www.govinfo.gov/content/pkg/USCOURTS-ca2-20-04238/pdf/USCOURTS-ca2-20-04238-0.pdf}{\textit{Melendez v. City of New York}, 16 F.4th 992 (2d Cir. 2021)}, the Second Circuit revived a Contracts Clause challenge to New York City's permanent cancellation of commercial-lease personal guaranties, and in \href{https://law.justia.com/cases/federal/appellate-courts/ca8/21-1278/21-1278-2022-04-05.html}{\textit{Heights Apartments, LLC v. Walz}, 30 F.4th 720 (8th Cir. 2022)}, the Eighth Circuit allowed a Contracts Clause claim against Minnesota's indefinite eviction moratorium to proceed past dismissal.
  \item \textbf{Post-Pandemic Deference Through 2026:} As \href{https://scholarship.law.wm.edu/cgi/viewcontent.cgi?filename=3&article=1014&context=propertyjournal&type=additional}{Professor James W. Ely, Jr. observed in 2025}, the Supreme Court has still not invalidated a state statute under the Contracts Clause in over forty years. Recent 2026 appellate decisions continue to apply \textit{Sveen} deferentially: in \href{https://law.justia.com/cases/federal/appellate-courts/ca2/25-1506/25-1506-2026-07-13.html}{\textit{Real Estate Board of New York, Inc. v. City of New York}, No. 25-1506 (2d Cir. July 13, 2026)}, the Second Circuit rejected a Contracts Clause preliminary-injunction challenge to New York City's broker-fee reform (the FARE Act), and in \href{https://law.justia.com/cases/new-york/appellate-division-first-department/2026/index-no-160266-20-appeal-no-5735-case-no-2025-00996.html}{\textit{513 W. 26th Realty LLC v. George Billis Galleries, Inc.}, 2026 NY Slip Op 06058 (N.Y. App. Div. 1st Dep't Oct. 20, 2026)}, New York's First Department upheld NYC's pandemic guaranty law under substantial legislative deference, expressly noting that the trial court erred by relying on Justice Gorsuch's \textit{Sveen} dissent rather than the controlling majority opinion.
\end{itemize}

\section{Pedagogical Critiques and Exam Hypotheticals}
\subsection{Pedagogical Critiques}
\begin{itemize}
  \item \textbf{Notice and Trap-for-the-Unwary Critique:} Justice Kagan's ``stroke of a pen'' rationale assumes policyholders know the legislature retroactively changed their pre-existing contract. An inattentive policyholder who relies on the law as it stood when contracting has no reason to file a new form to preserve an existing designation.
  \item \textbf{Doctrinal Collapse of Step One:} By folding probable intent and legislative reasonableness into step one (``substantial impairment''), \textit{Sveen} allows courts to dismiss Contracts Clause claims at the threshold without ever scrutinizing whether the State chose a narrowly tailored means at step two.
  \item \textbf{Structural Placement Note:} Although placed alongside federal legislative power and economic-regulation limits in this Briefing Book, remember that the Contracts Clause binds \textit{States} (Art. I, \S~10), not Congress---retroactive federal legislation impairing private contracts is instead evaluated under Fifth Amendment Due Process rationality review (or the Takings Clause).
\end{itemize}

\subsection{Exam Hypotheticals and ``Opposite Side'' Contrast}
\begin{itemize}
  \item \textbf{Hypothetical 1 (Irrevocable Rewriting of Private Contracts):} Suppose a State enacts a statute retroactively voiding all pre-existing non-compete clauses and personal loan guaranties in commercial leases executed before 2026, with no opt-out or reaffirmation mechanism. \textit{Analysis:} Unlike \textit{Sveen}, this law fails all three step-one factors---it contradicts party intent, shatters core reliance expectations, and permanently bars parties from safeguarding or reinstating their rights. The court must proceed to step two (\textit{Energy Reserves} / \textit{Melendez}) to determine whether the permanent abrogation is reasonable and appropriate to a significant public purpose.
  \item \textbf{Hypothetical 2 (``Opposite Side'' Contrast --- Public Contracts and Mandatory Paperwork vs. State Self-Interest):} Contrast \textit{Sveen} with \textit{United States Trust Co. v. New Jersey}, 431 U.S. 1 (1977), and \textit{Allied Structural Steel Co. v. Spannaus}, 438 U.S. 234 (1978). In \textit{Sveen}, the State acted as a neutral arbiter supplying a reversible default rule for private domestic-relations transfers. On the ``opposite side,'' when a State retroactively impairs its \textit{own} financial obligations (such as repealing a bond covenant or cutting vested public-employee pension formulas to save state revenue), courts withhold legislative deference at step two because the State's financial self-interest is directly at stake---demanding proof that the impairment was both reasonable and \textit{necessary} to serve an important public purpose.
\end{itemize}

\end{document}
